\section{Описание практической части}
\label{sec:Chapter4} \index{Chapter4}
Изучив базовый набор инструкций целевой виртуальной машины, перейдем непосредственно к реализации
верификаторных обработчиков инструкций. Причем верификация должна быть реализована
таким образом, чтобы выполнялись проверки на соответствие семантике, описанной в ISA.

Практическая часть работы была
сосредоточена на расширении возможностей верификатора байт-кода для ETS-специфичных инструкций и
инструкций взаимодействия с динамической средой. В техническом плане работа состояла из четырех
последовательных этапов. На первом этапе была подготовлена типовая инфраструктура для корректной
работы с рантайм-объединениями. На втором этапе была реализована верификация инструкций семейства
\texttt{ets.call.name.*}. На третьем этапе было добавлено покрытие для множества \texttt{any.*}.
Финальным этапом являлось тестирование реализованных обработчиков.

\subsection{Подготовка типовой инфраструктуры для рантайм-объединений}

Первым этапом был внесен подготовительный набор изменений, необходимый для более точного отображения
рантайм union в типовой системе верификатора. До этого в коде верификатора существовали места, где
union-тип фактически обрабатывался приближенно: представлял собой тип Object, часть проверок
именованных доступов к членам
объединений были ослаблены, поскольку информация о составе union не отображалась в
верификаторную модель достаточно точно.

В рамках этого этапа была переработана связь между \newline рантайм-представлением union и внутренним типом
верификатора. В частности, union-класс, существующий на уровне рантайма, перестал
рассматриваться как обычный class-тип. Для него была добавлена отдельная ветвь логики, которая
начала распознавать не только специальные union-значения верификатора, но и
рантайм-классы, помеченные как union. Это позволило извлекать множество составляющих типов
непосредственно из рантайм-описания класса и использовать их в операциях подтипизации и анализа
сигнатур.

\begin{lstlisting}[language=Python, label={lst:common_context_impl}, caption={Пример кода для поиска общего контекста}]
ClassLinkerContext *CoreClassLinkerExtension::GetCommonContext(Span<Class *> classes)
{
    ASSERT(!classes.Empty());
    auto *commonCtx = classes[0]->GetLoadContext();

    size_t foundClassesNum = 0;
    while (!commonCtx->IsBootContext()) {
        for (auto *klass : classes) {
            if (commonCtx->FindClass(klass->GetDescriptor()) != nullptr) {
                foundClassesNum++;
            }
        }
        if (foundClassesNum == classes.Size()) {
            break;
        }
        commonCtx = GetBootContext();
    }
    return commonCtx;
}
\end{lstlisting}

Дополнительно была доработана работа с контекстами загрузки классов. Для union-типа важно, чтобы
верификатор мог выбрать общий контекст, в котором корректно видны все составляющие классы
объединения. С этой целью в расширение class linker был добавлен алгоритм поиска общего
контекста загрузки для множества классов. Если общий контекст обнаружить не удается, логика
консервативно поднимается к boot context - это корневой контекст виртуальной машины, 
в котором находятся системные и глобально доступные классы платформы.

Корректность этого этапа была подтверждена отдельным verifier-тестом на сигнатуры union-типов.
Смысл теста заключался в проверке того, что сигнатура метода, принимающего объединение типов,
сопоставляется верификатором с рантайм-описанием union без перехода к излишне грубому общему
супертипу.

\subsection{Реализация верификации \texttt{ets.call.name.*}}

После подготовки типовой инфраструктуры была реализована верификация инструкций
\texttt{ets.call.name.short}, \texttt{ets.call.name} и \newline \texttt{ets.call.name.range}. Эти инструкции
используются для вызова метода по имени, в том числе в сценариях, связанных с объединениями типов.
До внесения изменений соответствующие обработчики в абстрактном интерпретаторе оставались
заглушками, поэтому верификатор не умел содержательно анализировать такой вызов.

Основное изменение было внесено в абстрактный интерпретатор. Для всех трех форм инструкции
была реализована единая схема анализа через общий вспомогательный интерфейс. В
рамках этой схемы верификатор:
\begin{enumerate}
    \item извлекает целевой метод из кэша, сформированного при разборе задания верификации;
    \item проверяет успешность разрешения \texttt{method\_id};
    \item выполняет проверку нарушений правил доступа к вызываемому методу;
    \item трактует первый аргумент как \texttt{this} и проверяет, что он является подтипом
    ссылочного типа;
    \item проверяет совместимость остальных фактических аргументов с формальной сигнатурой метода;
    \item устанавливает тип аккумулятора в тип возвращаемого значения вызываемого метода.
\end{enumerate}

Для таких форм, как \texttt{short}, работающей с двумя регистрами, и \texttt{full}, работающей с четырьмя регистрами,
набор аргументов извлекается из фиксированного числа
регистров. 
\begin{lstlisting}[language=Python, label={lst:range_impl}, caption={Пример кода для обработки \texttt{ets.call.name.range}}]
    template <BytecodeInstructionSafe::Format FORMAT>
    bool HandleEtsCallNameRange()
    {
        LOG_INST();
        DBGBRK();
        uint16_t vs = inst_.GetVReg<FORMAT, 0x00>();
        Method const *method = GetCachedMethod();
        if (method != nullptr) {
            LOG_VERIFIER_DEBUG_METHOD(method->GetFullName());
        }

        Sync();
        constexpr size_t INIT_SIZE = 5;
        SmallVector<int, INIT_SIZE> regs;
        for (auto regIdx = vs; ExecCtx().CurrentRegContext().IsRegDefined(regIdx); regIdx++) {
            regs.push_back(regIdx);
        }
        return CheckEtsCall<FORMAT>(method, Span {regs});
    }
\end{lstlisting}

Для \texttt{range}-формы был реализован отдельный сбор диапазона регистров из текущего
верификаторного контекста, см. листинг ~\ref{lst:range_impl}, что согласует проверку с соглашением о кодировании аргументов в ISA.

С практической точки зрения важным было решение не дублировать всю рантайм-семантику инструкции в
верификаторе. Проверка была построена как консервативная статическая аппроксимация: на этапе
верификации проверяются только те свойства, которые могут быть надежно установлены без исполнения
программы. В частности, верификатор гарантирует, что первый аргумент, интерпретируемый как
\texttt{this}, имеет ссылочный тип, что остальные фактические аргументы совместимы с формальной
сигнатурой вызываемого метода, и что тип результата, помещаемого в аккумулятор, согласован с
объявленным типом возвращаемого значения. На этапе исполнения остается позднее разрешение целевого 
метода по фактическому классу объекта: после извлечения
описания метода требуется подобрать совместимую реализацию именно для того класса,
который находится в регистре в момент исполнения. Данный сценарий особенно важен в контексте объединения типов,
где статически известен лишь абстрактный тип, а конкретный класс получателя
определяется только в рантайме. Также в интерпретатор были вынесены проверки на отсутствие подходящего 
метода в фактическом классе и на случай, когда найденная реализация остается абстрактной.

Для подтверждения корректности были добавлены как unit-тесты, так и негативные ассемблерные программы.
Положительные тесты покрывают все три формы инструкции: короткую, полную и диапазонную.

Таким образом, на этом этапе верификатор получил поддержку ETS-специфичного
множества инструкций, ранее не покрытого статическим анализом.

\subsection{Реализация верификации \texttt{any.*}}

Третьим этапом была реализована поддержка семейства \texttt{any.*}, отвечающего за операции над
значениями, проходящими через интерфейс взаимодействия со статической и динамической частями
платформы. В ходе работы были доработаны правила в ISA,
расширен набор проверок в абстрактном интерпретаторе и добавлен набор тестов.

Первое изменение коснулось самой спецификации верификации в \newline \texttt{isa.yaml}. Для инструкций
\texttt{any.ldbyidx} и \texttt{any.stbyidx} были уточнены типы индексных операндов: вместо
\texttt{i32} в правилах проверки были зафиксированы \texttt{f64}-значения, что соответствует
реальной семантике инструкции. Кроме того, для \texttt{any.isinstance} был уточнен тип результата в
аккумуляторе: результат проверки должен интерпретироваться как \texttt{i32}.

В абстрактном интерпретаторе были добавлены две служебные функции:
\texttt{CheckRegTypeNotNullRef()} и \texttt{CheckRegTypes()}. Первая выделяет типовую для этого
семейства проверку случая, когда инструкция гарантированно приведет к \texttt{NullPointerException}.
Вторая позволяет компактно проверять несколько регистров сразу, что существенно упростило код
обработчиков.

После этого были реализованы обработчики следующих инструкций:
\begin{itemize}
    \item вызовы функций и методов:
    \texttt{any.call.0}, \texttt{any.call.short}, \newline \texttt{any.call.range},
    \texttt{any.call.this.0}, \texttt{any.call.this.short}, \newline \texttt{any.call.this.range};
    \item вызовы конструкторов:
    \texttt{any.call.new.0}, \texttt{any.call.new.short}, \newline \texttt{any.call.new.range};
    \item доступ к свойствам и индексаторам:
    \texttt{any.ldbyname}, \texttt{any.ldbyname.v}, \texttt{any.stbyname},
    \texttt{any.stbyname.v}, \texttt{any.ldbyidx}, \texttt{any.stbyidx},
    \texttt{any.ldbyval}, \texttt{any.stbyval};
    \item проверка типа:
    \texttt{any.isinstance}.
\end{itemize}

Для большинства этих инструкций структура верификации оказалась однотипной. Сначала выполняется
синхронизация состояния абстрактной интерпретации. Затем проверяются типы
операндов: ссылочные аргументы должны иметь ссылочный тип, индексные аргументы --- тип
\texttt{f64}, а аккумулятор должен содержать значение ожидаемого типа. После этого выполняется
отдельная проверка на случай, когда значение объекта-приемника заведомо является нулевым. Если такая ситуация 
обнаруживается, обработчик завершает анализ с фиксацией случая \texttt{AlwaysNpe}. 
Иначе верификатор обновляет тип результата:
либо задает тип аккумулятора, либо выполняет переход к следующей
инструкции.

Стоит отметить, что реализация была построена максимально единообразно.
Семейство \texttt{any.*} содержит много инструкций с близкой семантикой, поэтому выделение общих
проверок позволило избежать большого количества повторяющегося кода.

\subsection{Тестирование результатов}

Для инструкций \texttt{any.*} был добавлен отдельный набор unit-тестов. 
Этот набор покрывает как
положительные, так и отрицательные сценарии. В него вошли тесты на вызовы функций и методов,
конструкторы, чтение и запись по имени, чтение и запись по индексу, доступ по значению-ключу и
инструкцию \texttt{any.isinstance}. Отдельные негативные тесты проверяют несовместимость типов,
например использование неверного типа индексного операнда.

Также в тестовую инфраструктуру была внесена важная доработка: была реализована
удобная инфраструктура для тестирования новых верификаторных обработчиков инструкций, которая упростила
процесс проверки и обеспечила корректное выполнение verifier-тестов.

\subsection{Итоги практической части}

Итогом практической части стало расширение множества инструкций, для которых верификатор способен
выполнять содержательную проверку типов и состояний регистров. Работа была построена поэтапно: от
подготовки типовой модели объединенных типов, через добавление поддержки ETS-специфичных
вызовов, к реализации крупного блока проверок для \texttt{any.*}. Такой порядок оказался
принципиально важным, поскольку последующие изменения опирались на корректное отображение рантайм-
типов в модель верификатора.

С практической точки зрения проделанная работа улучшила две ключевые характеристики системы.
\begin{itemize}
    \item \textbf{Безопасность:} был увеличен класс верифицируемых статически программ.
    \item \textbf{Совместимость:} повысилась точность верификации за счет более тесного соответствия между рантайм-семантикой
платформы и внутренней типовой системой верификатора.
\end{itemize}
